% ==============================================================================
% 3. Feladat: Komponensválasztás és Numerikus Szimuláció
% ==============================================================================

\section{Feladat: Komponensválasztás és Numerikus Szimuláció}

\begin{problem}[15 pont]
Válasszon megfelelő képérzékelő szenzort és ipari optikát a specifikáció kielégítésére! Foglalja össze a paramétereket a Style B (2 oszlopos, multirow) táblázatformátumban, majd mellékelje a szabályozási kört szimuláló MATLAB forráskódot!
\end{problem}

\subsection{Kiválasztott optikai és érzékelő komponensek (Style B Táblázat)}
A választott képérzékelő modul és ipari optika főbb specifikációit a \textbf{2. táblázat} tartalmazza a gyári adatlap alapján \cite{sony_imx_datasheet}.

% --- STYLE B: 2-Column Component Specification Table with Multirow ---
\begin{table}[ht!]
\centering
\renewcommand{\arraystretch}{1.25}
\setlength{\tabcolsep}{8pt}
\setlength{\arrayrulewidth}{2pt}
\arrayrulecolor{white}
\begin{tabularx}{\textwidth}{|X!{\color{white}\vrule width 3pt}X|}
\hline
\rowcolor{tableGray}
\textbf{Optikai és Szenzor Paraméter} & \textbf{Specifikált Érték / Típus} \\ \hline
\rowcolor{white}
Képérzékelő típusa & Sony IMX477 CMOS (Type 1/2.3) \\ \hline
\rowcolor{tableGray}
Aktív felbontás & $4056 \times 3040\text{ (12{,}3 MP)}$ \\ \hline
\rowcolor{white}
 & Vízszintes: $84{,}7^\circ$ \\
\rowcolor{white}
 & Függőleges: $66{,}1^\circ$ \\
\rowcolor{white}
\multirow{-3}{\linewidth}{Látómező (FOV a szenzorméretre)} & Diagonális: $100{,}1^\circ$ \\ \hline
\rowcolor{tableGray}
Objektív gyújtótávolsága ($f$) & $4{,}50\text{ mm}$ (Fix fókuszú) \\ \hline
\rowcolor{white}
Tárgytávolsági tartomány ($s$) & $25\text{ mm} - \infty$ \\ \hline
\rowcolor{tableGray}
Optikai szűrő tartománya & $400 - 700\text{ nm}$ (VIS látható spektrum) \\ \hline
\rowcolor{white}
Interfész és kommunikáció & MIPI CSI-2 (2 vagy 4 sávos) \\ \hline
\end{tabularx}
\caption{A választott ipari optikai és detektormodul műszaki specifikációi \cite{sony_imx_datasheet}}
\label{tab:spec_tabla}
\end{table}

\subsection{Szimulációs MATLAB szkript}
A PID szabályozó és a mechanikai szakasz zárt köri vizsgálatát az \textbf{1. forráskód} szerinti MATLAB szkript végzi.

\begin{lstlisting}[language=Matlab, caption={PID szabályozási kör szimulációja és Bode-analízis}]
% =========================================================================
% BME MOGI - Szabalyozastechnika es Muszaki Optika szimulacio
% Hallgato: Minta Janos (NEPTUN)
% =========================================================================
clear; clc; close all;

% Szakasz parameterei (1. feladat alapjan)
J = 2.5e-3;   % Redukalt tehetetlensegi nyomatek [kg*m^2]
b = 1.2e-2;   % Csillapitasi tenyezo [N*m*s/rad]
k = 0.85;     % Visszaterito rugomerevseg [N*m/rad]
Kt = 0.45;    % Motor nyomatektenyezo [N*m/A]

s = tf('s');
P = Kt / (J*s^2 + b*s + k);  % Szakasz atviteli fuggveny

% PID szabalyozo tervezes
Kp = 14.5;
Ti = 0.22;
Td = 0.045;
C = Kp * (1 + 1/(Ti*s) + Td*s / (1 + 0.005*s));

% Zart kor es vizsgalatok
L = C * P;              % Felnyitott kor
T = feedback(L, 1);     % Zart kor alapjelre
S = feedback(1, L);     % Erzekenysegi fuggveny

figure('Name', 'Zart kor lepesvalasza');
step(T, 10);
grid on;
title('BME MOGI - Pozicionalo kor lepesvalasza');
xlabel('Ido [s]'); ylabel('Szogelfordulas [rad]');
\end{lstlisting}
